
\section{\trBertini{}定理}

现在我们来讨论\trBertini{}定理——一条基本的经典结果。

\textbf{13.4.1.}\ \emph{\trhDefinition{}：\trDualProjectiveSpace。}\quad $\mathbb{P}^n_k$（具有射影坐标 $x_0,\dots,x_n$）的\textbf{\trDual}（dual）（或\textbf{\trDualProjectiveSpace}（dual projective space）），非正式地说，就是 $\mathbb{P}^n_k$ 中\trHyperplane{}所成的空间。说得稍稍确切些：它是一个具有射影坐标 $a_0,\dots,a_n$ 的射影空间 $\mathbb{P}^n_k$（我们把它记作 $\mathbb{P}^{n\vee}_k$，徒劳地指望借此防止混淆），并附带如下数据：「\trIncidenceVariety」或「\trIncidenceCorrespondence」$I\subset\mathbb{P}^n\times\mathbb{P}^{n\vee}$，它由方程 $a_0x_0+\cdots+a_nx_n=0$ 切出。（这常称为\textbf{\trUniversalHyperplane}（universal hyperplane）。）注意：$\mathbb{P}^{n\vee}$ 的 $k$-值点确实对应于 $\mathbb{P}^n$ 中定义在 $k$ 上的那些\trHyperplane，而且这显然还是一种\trDuality{}关系（在定义中，$x$-变量与 $a$-变量之间存在着对称性）。这个定义足够具体，可供实践使用，并且能延拓到任意的\trBase（特别是 $\Spec\mathbb{Z}$）之上。（但如果你品味精致细腻，或许会想给出一个\trCoordinateFree{}的定义。）

\textbf{13.4.2.\ \trBertini{}定理。}\ ——\ \emph{设 $X$ 是 $\mathbb{P}^n_k$ 中（纯）维数为 $d$ 的一个\trSmooth\trClosed\trSubvariety。那么在\trDualProjectiveSpace $\mathbb{P}^{n\vee}_k$ 中存在一个非空（＝稠密）的\trOpenSubset{} $\mathcal{U}$，使得对每一点 $p=[H]\in\mathcal{U}$，概形 $H\cap X$ 在 $\kappa(p)$ 上光滑，且具有（纯）维数 $d-1$。}

（你或许应当先想想 $\kappa(p)=k$ 的情形，即 $H$ 是由一条系数取自 $k$ 的线性型所切出的\trHyperplane。）凡属这种风味的定理，常统称为「\trBertini{}定理」。一个例子是\trKleiman{}--\trBertini{}定理~21.6.10，它可\emph{不是}由\trKleiman{}与\trBertini{}二人联手证明的。

作为\trBertini{}定理~13.4.2 的一个应用：$\mathbb{P}^n_k$ 中一张一般的次数 $d>0$ 的超曲面与 $X$ 交出 $X$ 中\trCodimension{}为 $1$ 的一个光滑\trSubvariety：把 $X\hookrightarrow\mathbb{P}^n$ 换成复合
\[
\begin{tikzcd}
X \ar[r] & \mathbb{P}^n \ar[r, hook, "\nu_d"] & \mathbb{P}^N,
\end{tikzcd}
\]
其中 $\nu_d$ 是第 $d$ 次\trVeroneseEmbedding{}（9.3.8）。这里「一般」取其在\trAlgebraicGeometry{}中的惯常含义；见 \S10.3.7。（\mbox{练习~13.4.B}给出一个推广。）这样做的一个有用推论是：取 $X=\mathbb{P}^n$，我们立即看到 $\mathbb{P}^n$ 中存在光滑的次数 $d$ 超曲面（在任何无穷域 $k$ 上），而无须同具体的方程纠缠（\mbox{练习~13.3.E}即其特殊情形）。

\mbox{练习~21.6.I}给出了特征 $0$ 时\trBertini{}定理的一个有用的改进（见\mbox{练习~21.6.J}）。

\emph{证明。}\ 这条证明的中心想法相当朴素。我们将把那些「坏的」超平面描述出来，并证明它们构成 $\mathbb{P}^{n\vee}_k$ 的一个维数至多为 $n-1$ 的闭子集，从而其补集是一个稠密开子集。说得稍稍确切些：我们将定义一个射影簇 $Z\subset X\times\mathbb{P}^{n\vee}_k$，它可以非正式地描述为
\begin{equation*}\tag{13.4.2.1}\label{eq:bertini-bad-locus-z}
\begin{split}
Z=\bigl\{(p\in X,\ H\subset\mathbb{P}^n_k)&\ :\ p\in H,\\
&\text{并且 }p\text{ 是 }H\cap X\text{ 的\trSingularPoint，或者 }X\subset H\bigr\}.
\end{split}
\end{equation*}
我们将见到：投射 $\pi\colon Z\to X$ 在闭点处的纤维是维数为 $n-1-\dim X$ 的射影空间，从而 $\dim Z\le n-1$。于是 $Z$ 在 $\mathbb{P}^{n\vee}$ 中的像将是（定理~8.4.10）一个维数至多为 $n-1$ 的闭子集，故其补集既开又非空。现在我们就把这一策略付诸实施。

\textbf{13.4.3.}\quad 我们首先更精确地定义 $Z$，即通过 $\mathbb{P}^n\times\mathbb{P}^{n\vee}$ 上的方程来定义：$\mathbb{P}^n$ 上的坐标是 $x_0,\dots,x_n$，而 $\mathbb{P}^{n\vee}$ 上的对偶坐标是 $a_0,\dots,a_n$。设 $X$ 由 $f_1,\dots,f_r$ 切出。那么我们就取这些方程作为 $Z$ 的定义方程中的第一批。（至此我们定义了子概形 $X\times\mathbb{P}^{n\vee}$。）我们再添加方程 $a_0x_0+\cdots+a_nx_n=0$。（至此我们描述了 $\mathbb{P}^n\times\mathbb{P}^{n\vee}$ 中由点 $(p,H)$——其中 $p\in X$ 且 $p\in H$——所对应的子概形。）注意：射影\trJacobianMatrix{} \eqref{eq:projective-jacobian-corank}
\[
\begin{pmatrix}
\dfrac{\partial f_1}{\partial x_0}(p) & \cdots & \dfrac{\partial f_r}{\partial x_0}(p)\\
\vdots & \ddots & \vdots\\
\dfrac{\partial f_1}{\partial x_n}(p) & \cdots & \dfrac{\partial f_r}{\partial x_n}(p)
\end{pmatrix}
\]
由射影雅可比\trCorank{}判别法（\mbox{练习~13.1.N}）知，在 $X$ 的所有点处的\trCorank{}恰为 $d+1$。于是我们要求：在射影\trJacobianMatrix{}上添加新的一列
\[
\begin{pmatrix} a_0\\ \vdots\\ a_n \end{pmatrix}
\]
之后，所得矩阵的\trCorank{}不小于 $d+1$（因而 $=d+1$）。这恰是我们要捕捉的观念——你明白这是为什么吗？它由若干方程切出（某些\trMinor{}的行列式）。（注记~13.4.4 完整地走了一遍例子，或许有助于澄清这里的运作方式。）

我们接着证明 $\dim Z\le n-1$。对每个\trClosedPoint{} $p\in X$，令 $W_p$ 为包含 $p$、且满足「$H\cap X$ 在 $p$ 处\trSingular，或 $H$ 包含整个 $X$」的那些超平面所成的\trLocus；$W_p$ 的维数是多少？设 $\dim X=d$。那么 $W_p$ 的定义中对超平面的限制相应于 $d+1$ 个线性条件。（你明白为什么吗？）这意味着 $W_p$ 是一个\trCodimension{}为 $d+1$（也就是维数为 $n-d-1$）的射影空间。于是 $\pi\colon Z\to X$ 在每个\trClosedPoint{}上的纤维都具有纯维数 $n-d-1$。由\mbox{关键练习~12.4.A}，这蕴含 $\dim Z\le n-1$。（如果你愿意，可以用\mbox{练习~12.4.D}证明 $\dim Z=n-1$；并且日后等你知道了什么是射影丛（定义~17.2.4），还可以证明 $Z$ 是 $X$ 上的一个射影丛。但本证明用不到这些。）\hfill$\square$

\textbf{13.4.4.}\ \emph{\trhRemark{}。}\quad 这里有一个例子，或许有助于让你确信：段落~13.4.3 中的代数确实描述了我们想要的几何。考虑特征不为 $2$ 的域上的平面\trConic{} $x_0^2-x_1^2-x_2^2=0$；从坐标卡 $U_0$ 中的实图像着眼，你不妨把它想象成圆 $x^2+y^2=1$。考虑点 $[1,1,0]$，它对应于圆上的 $(1,0)$。我们期待仿射平面中的\trTangentLine{}是 $x=1$，它相应于 $x_0-x_1=0$。让我们看看代数给了我们什么。射影\trJacobianMatrix{}为
\[
\begin{pmatrix} 2x_0\\ -2x_1\\ -2x_2 \end{pmatrix}
=\begin{pmatrix} 2\\ -2\\ 0 \end{pmatrix},
\]
它果然如预期般具有秩 $1$。我们的配方要求矩阵
\[
\begin{pmatrix} 2 & a_0\\ -2 & a_1\\ 0 & a_2 \end{pmatrix}
\]
具有秩 $1$（即 $a_0=-a_1$ 且 $a_2=0$），并且还要求 $a_0x_0+a_1x_1+a_2x_2=0$，而这（请你自行核验）恰是我们想要的。

\textbf{13.4.A.（容易的练习）}
\begin{enumerate}[label=(\alph*)]
\item 在较弱的假设「$X$ 只有有限多个非光滑点」之下证明\trBertini{}定理。
\item 在较弱的假设「$X\to\mathbb{P}^n_k$ 是\trLocallyClosedEmbedding（但 $X$ 仍光滑）」之下证明\trBertini{}定理。
\end{enumerate}

\textbf{13.4.B.（练习：多超曲面的\trBertini{}定理）}\quad 设 $X$ 是 $\mathbb{P}^n_k$ 中（纯）维数为 $d\ge r$ 的一个光滑\trSubvariety。设 $e_1,\dots,e_r$ 是正整数，于是
\begin{equation*}\tag{13.4.4.1}\label{eq:multi-hypersurfaces-param}
\mathbb{P}^{\binom{e_1+n}{n}-1}\times\cdots\times\mathbb{P}^{\binom{e_r+n}{n}-1}
\end{equation*}
参数化了次数分别为 $e_1,\dots,e_r$ 的超曲面所成的 $r$ 元组 $(H_1,\dots,H_r)$。证明：存在 \eqref{eq:multi-hypersurfaces-param} 的一个稠密开子集 $\mathcal{U}$，使得对每一点 $p=[(H_1,\dots,H_r)]\in\mathcal{U}$，$H_1\cap\cdots\cap H_r\cap X$ 在 $\kappa(p)$ 上光滑，且具有（纯）维数 $d-r$。非正式地说：对次数分别为 $e_1,\dots,e_r$ 的超曲面 $H_1,\dots,H_r$ 作一般选取，则 $H_1\cap\cdots\cap H_r\cap X$ 光滑，且具有（纯）维数 $d-r$。（提示：不要去套用\trBertini{}定理~13.4.2，而是试着给出「同样的」证明。）这道练习的要点不在于要你记住这个具体结果，而在于：凡你需要的那类结果，你都有能力亲手证明。

\textbf{13.4.C.（练习）}\quad 设 $X$ 是 $\mathbb{P}^n$ 中维数为 $d$ 的一个射影簇，带有一个光滑的稠密开子集。（若基域完全，这件事会自动成立：见关于簇的一般光滑性（generic smoothness）的定理~21.6.1。）证明：$X$ 与 $d$ 张一般\trHyperplane{}的交由有限多个既约点组成。更精确地说：若 $\mathbb{P}^{n\vee}$ 是\trDualProjectiveSpace，则存在一个非空的\trZariski{}开子集 $\mathcal{U}\subset(\mathbb{P}^{n\vee})^d$，使得对 $\mathcal{U}$ 的每个\trClosedPoint{} $(H_1,\dots,H_d)$，概形论交 $H_1\cap\cdots\cap H_d\cap X$ 由有限多个既约点组成。（这样的点的个数——按正确方式计数——称为该簇的\trDegree；见\mbox{练习~18.6.M}。）

\textbf{13.4.5.}\ \emph{\trDualVariety。}\quad 我们关于\trBertini{}定理的讨论，给了我们一个顺便提及一个经典构造的借口。$Z$（见 \eqref{eq:bertini-bad-locus-z}）在 $\mathbb{P}^{n\vee}$ 中的像称为 $X$ 的\textbf{\trDualVariety}（dual variety）。由于 $\dim Z=n-1$，我们「期待」$X$ 的对偶是 $\mathbb{P}^{n\vee}$ 中的一张超曲面。这是一个并不显然的事实：这确实是一种\trDuality——$X$ 的对偶的对偶就是 $X$ 自身（见 [H2, Thm.~15.24]）。下面这道练习会让你对\trDualVariety{}有几分感觉。

\textbf{13.4.D.（练习）}\quad 证明：$\mathbb{P}^n$ 中一张\trHyperplane{}的对偶，就是对偶空间 $\mathbb{P}^{n\vee}$ 中那个相应的点。如此一来，$\mathbb{P}^n$ 与 $\mathbb{P}^{n\vee}$ 之间的对偶，便是射影簇之间对偶的一个特例。

\textbf{13.4.E.（练习）}\quad 设 $C\subset\mathbb{P}^2$ 是特征不为 $2$ 的\trAlgClosedField{}上的一条光滑\trConic。证明：$C$ 的\trDualVariety{}也是一条光滑\trConic。于是，举例来说，过平面上一个一般点（若 $k=\overline{k}$），有 $C$ 的两条切线。（对偶平面中一条直线上的诸点，对应于原平面中过某一点的那些直线。）

\textbf{13.4.F.}\ \hard\ \textbf{（练习：恰有一条光滑\trConic{}与五条一般直线相切，及其推广）}\quad 沿用上一题的记号，证明：含 $i$ 个一般选定的点、且与 $5-i$ 条一般选定的直线相切的光滑\trConic{} $C$ 的个数，当 $i=0$、$1$、$2$、$3$、$4$、$5$ 时，分别为 $1$、$2$、$4$、$4$、$2$、$1$。你不妨借助\trConic{}与对偶\trConic{}之间的对偶，来解读这个数列的对称性。这一事实恐怕早在旧石器时代就已为人所知了。
